\ifdefined\mainfile\else
\documentclass{article}
\usepackage{graphicx}
\usepackage{float}
\usepackage{amsmath}
\usepackage{amssymb}
\usepackage[a4paper, left=2.5cm, right=2.5cm, top=2.5cm, bottom=2.5cm]{geometry}
\usepackage[colorlinks=true,linkcolor=black,citecolor=blue,urlcolor=blue]{hyperref}
\begin{document}
\fi

\section{Datapath}
\label{sec:datapath}
% Top-level wiring: RegisterFile -> MUXB -> FunctionUnit -> MUXD -> write-back

\subsection{MUXB}
\label{sec:muxb}
MUXB er en 8-bit 2-to-1 multiplexer, der vælger hvilken operand der føres ind som B-input til Function Unit. Multiplexeren styres af kontrolsignalet MB, som bestemmer om data fra Register File eller en ekstern konstant bruges som operand.
Når MB er 0, vælges B\_Data fra Register File, som er den værdi der læses fra det register som er valgt af B-adressen. Når MB er 1, vælges ConstantIn, hvilket gør det muligt at bruge en immediate konstant direkte i en mikrooperation.
MUXB er implementeret som en instans af MUX2x1x8 komponenten, hvor J er forbundet til B\_Data, H til ConstantIn, og MF til MB signalet. Outputtet Bus\_B føres videre som B-operand til Function Unit.


\subsection{MUXD}
\label{sec:muxd}
MUXD er en 8-bit 2-to-1 multiplexer, der vælger hvilken data der skrives tilbage til Register File som D\_Data. Multiplexeren styres af kontrolsignalet MD, som bestemmer om resultatet fra Function Unit (F) eller ekstern data (DataIn) føres videre.
Når MD er 0, vælges resultatet F fra Function Unit, hvilket bruges ved aritmetiske og logiske operationer. Når MD er 1, vælges DataIn, som gør det muligt at loade ekstern data direkte ind i registrene. MUXD er implementeret som en instans af MUX2x1x8 komponenten, hvor J er forbundet til F\_Out, H til DataIn, og MF til MD signalet.

\subsection{Datapath Test}
\label{sec:datapath-test}
% Integrations-testbench -- 19 tests

Vi har lavet en testbench, som tester at Datapath fungerer korrekt. 
\textbf{Fig \ref{fig:Datapath_sim}} viser simuleringen af den samlede Datapath. 
Simuleringen tester, at Datapath kan udføre alle funktionerne fra 
funktionstabellerne i opgavebeskrivelsen og bekræfter dermed, at vores 
implementering virker som forventet.

\begin{figure}[H]
    \centering
    \includegraphics[width=0.75\linewidth]{images//datapath/Datapath_sim.png}
    \caption{Simulering af Datapath}
    \label{fig:Datapath_sim}
\end{figure}

\ifdefined\mainfile\else
\end{document}
\fi
